\chapter{问题三模型与结果：异常识别和变量贡献}

\section{变量分类与补齐}

附件3同时包含降雨量、孔隙水压力、微震事件数、深部位移和表面位移。若对全部变量使用同一异常规则，会把真实降雨或微震事件误判为异常。因此本文先按物理性质分类。

孔隙水压力、深部位移和表面位移属于连续状态量，采用局部线性状态空间模型进行缺失补齐：
\begin{equation}
  z_t=l_t+\epsilon_t,\quad
  l_t=l_{t-1}+b_{t-1}+\eta_t,\quad
  b_t=b_{t-1}+\zeta_t .
\end{equation}
其中$z_t$为观测值，$l_t$为平滑状态，$b_t$为局部趋势。降雨量和微震事件数属于稀疏驱动量，大量零值是真实状态，只对极端强度和与位移残差同步的异常事件标记。补齐后数据无缺失，且降雨、微震等非负变量不出现负值。

\section{异常点结果}

单变量异常数量见表\ref{tab:q3_single_anomaly}。表面位移异常点最多，说明位移观测受短时扰动和阶段变化影响更强；降雨量异常较多但主要反映集中降雨事件，不能简单视为坏点。共同异常定义为同一编号处至少两个变量同时异常，结果见表\ref{tab:q3_common_anomaly}。

\begin{table}[htbp]
\centering
\small
\caption{单变量异常点统计（题目表3.1）}
\label{tab:q3_single_anomaly}
\begin{tabular}{lr}
\toprule
数据集变量 & 异常点数量 \\
\midrule
a：降雨量 & 17 \\
b：孔隙水压力 & 5 \\
c：微震事件数 & 0 \\
d：深部位移 & 5 \\
e：表面位移 & 36 \\
总数 & 63 \\
\bottomrule
\end{tabular}
\end{table}

\begin{table}[htbp]
\centering
\small
\caption{共同异常点（题目表3.2）}
\label{tab:q3_common_anomaly}
\begin{tabular}{rrrrrrrr}
\toprule
序号 & 编号 & 变量 & 序号 & 编号 & 变量 & 序号 & 编号/变量 \\
\midrule
1 & 1406 & ae & 7 & 3352 & ae & 13 & 8948/ae \\
2 & 1444 & ae & 8 & 4865 & ae & 14 & 9035/ae \\
3 & 1632 & ae & 9 & 5542 & ae & 15 & 9419/bd \\
4 & 1759 & de & 10 & 6528 & ae & 16 & 9491/de \\
5 & 2021 & de & 11 & 8105 & de & 17 & 9800/ae \\
6 & 3031 & ae & 12 & 8253 & ae &  &  \\
\bottomrule
\end{tabular}
\end{table}

异常变量组合主要有两类：一类是降雨量与表面位移同步异常，即$a$--$e$共同异常；另一类是深部位移与表面位移同步异常，即$d$--$e$共同异常。前者说明降雨入渗对浅表响应影响明显，后者说明深部剪切变形会传导到表面位移。

\begin{figure}[htbp]
\centering
\includegraphics[width=0.94\textwidth]{q3_variable_trace_panels.png}
\caption{问题三原始、补齐与异常点叠加图}
\label{fig:q3_trace}
\end{figure}

\begin{figure}[htbp]
\centering
\includegraphics[width=0.82\textwidth]{q3_missing_heatmap.png}
\caption{问题三缺失分布热力图}
\label{fig:q3_missing}
\end{figure}

图\ref{fig:q3_missing}用于检查缺失是否集中出现在某些变量或时段。若缺失高度集中，则简单线性插值会低估不确定性；本文采用状态空间平滑，正是为了在保留局部趋势的同时避免对长缺失段过度拟合。

\section{表面位移预测模型对照}

对补齐后的训练集构造滞后、滚动和交互特征，以表面位移为响应变量。为兼顾线性可解释性与非线性拟合能力，本文比较ElasticNet、GBDT和HistGBDT。GBDT类模型源于梯度提升树思想\citep{friedman2001greedy}，适合捕捉阈值型降雨、孔压和微震响应。模型对照结果见表\ref{tab:q3_model_compare}。

\begin{table}[htbp]
\centering
\small
\caption{问题三预测模型对照}
\label{tab:q3_model_compare}
\begin{tabular}{lrrr}
\toprule
模型 & RMSE/mm & MAE/mm & $R^2$ \\
\midrule
HistGBDT & 4.668 & 2.729 & 0.877 \\
GBDT & 4.728 & 2.763 & 0.874 \\
ElasticNet & 4.969 & 2.831 & 0.861 \\
\bottomrule
\end{tabular}
\end{table}

HistGBDT的验证误差最低，但GBDT与其差距很小，且能直接输出特征重要性并便于做置换贡献分析。因此本文把HistGBDT作为泛化性能对照，把GBDT作为后续贡献排序和实验集估计的主模型。

\section{异常阈值敏感性}

连续状态量的Hampel阈值从3.5增至5.5时，异常数量平滑下降，说明异常识别不是由某个临界参数偶然决定。表\ref{tab:q3_sensitivity}列出主要结果。本文最终取4.5作为连续变量阈值，并对降雨采用0.997分位数稀疏事件阈值，对应17个降雨异常点。

\begin{table}[htbp]
\centering
\small
\caption{问题三异常阈值敏感性}
\label{tab:q3_sensitivity}
\begin{tabular}{rrrrr}
\toprule
连续阈值 & 孔压异常 & 深部位移异常 & 表面位移异常 & 降雨异常(0.997分位) \\
\midrule
3.5 & 37 & 29 & 124 & 17 \\
4.0 & 16 & 10 & 68 & 17 \\
4.5 & 5 & 5 & 36 & 17 \\
5.0 & 0 & 4 & 21 & 17 \\
5.5 & 0 & 0 & 11 & 17 \\
\bottomrule
\end{tabular}
\end{table}

\section{变量贡献与实验集估计}

变量贡献采用两种口径：置换重要性衡量扰乱某类变量后RMSE的增加，重复划分贡献衡量稳定性。结果见表\ref{tab:q3_contribution}。孔隙水压力和深部位移贡献最高且标准差较小，是边坡位移响应的主控因子；降雨量主要通过改变孔压和短期残差发挥作用；微震事件数贡献较小，但在快速阶段仍有预警意义。

\begin{table}[htbp]
\centering
\small
\caption{问题三变量贡献稳定性}
\label{tab:q3_contribution}
\begin{tabular}{lrrr}
\toprule
变量组 & 置换RMSE增量/mm & 稳定性均值 & 稳定性标准差 \\
\midrule
孔隙水压力 & 10.071 & 9.786 & 0.427 \\
深部位移 & 9.244 & 8.811 & 0.192 \\
降雨量 & 3.779 & 3.348 & 0.347 \\
微震事件数 & 1.276 & 1.088 & 0.101 \\
\bottomrule
\end{tabular}
\end{table}

将训练完成的GBDT融合模型应用于实验集，得到5000个实验编号的表面位移估计值。预测结果范围为67.546--148.123 mm，中位数为93.808 mm。散点图见图\ref{fig:q3_exp_scatter}，其作用是给出题目要求的实验集位移估计结果，并检查预测值是否出现明显离群或非物理跳变。

\begin{figure}[htbp]
\centering
\includegraphics[width=0.82\textwidth]{q3_experiment_prediction_scatter.png}
\caption{问题三实验集表面位移估计散点图}
\label{fig:q3_exp_scatter}
\end{figure}
